\chapter[熵增是宇宙基本法则]{熵增是宇宙基本法则}

\chaptersubtitle{《我真是个笨蛋》}

\chapternote{%
    ``Just as \textbf{the constant increase of entropy is the basic law of the universe}, so it is the basic law of life to be ever more highly structured and to struggle against entropy.''
    \par
    {\raggedleft\noindent ---Václav Havel, \emph{Letter to Husák}\par}
    \par
    “正如\textbf{熵增是宇宙基本法则}，生命的基本法则也是不断形成更高层次的组织结构，并持续抵抗熵所代表的无序化趋势。”
    \par
    {\raggedleft\noindent ——瓦茨拉夫·哈维尔《致胡萨克的公开信》\par}
}

《魔法少女小圆》第八集呈现了这类作品中常见的一种叙事问题。
对于《魔圆》这样主要依赖强烈的剧情反转\footnote{%
    这种风格在日本很常见，但并非独有。
    它有着悠久的历史，常以一种名为“起承转合”的修辞和叙事结构呈现。
    该结构起源于中国四行诗，其中第一句“起”用于引入主题，第二句“承”进一步展开主题，第三句“转”引入一个惊人的转折或表面上无关的内容，第四句“合”将转折整合到原话题中，从而得出一个新颖或出乎意料的结论。
}来推动剧情的作品而言，突发的重大转折固然能够制造冲击力，却也很容易压过此前细腻的角色塑造\cite{maynard1997communication-strategies}。
因此，本章暂不重点讨论结尾的转折，而是先集中分析前十八分钟左右的角色叙事与情感铺陈。

动画开头的十八分钟可以说是全剧叙事焦点最集中的段落之一。
除去极少数场景，观众几乎始终与沙耶香在一起。
即便有一个场景表面上发生在恭介与仁美之间，似乎暂时脱离了沙耶香，后续仍然通过她的视角被重新组织起来。
因此，这一集几乎全部篇幅都用来集中呈现沙耶香走向最终崩溃的过程：
魔法少女身份带来的持续压力，丘比揭示身体真相后她自我价值感的彻底丧失，以及她拒绝净化灵魂宝石的行为，彻底摧毁了她的精神状态。

这一分析不仅解释了沙耶香的变化过程，也进一步揭示了“魔法少女”在该作品世界观中的核心含义之一：
魔法少女的存在，本质上就是持续与抑郁相抗衡。
作品对抑郁的描写相当贴近这种心理障碍的真实体验。
少女们每次使用魔法时，灵魂宝石中都会积累黑暗；
这种黑暗并不是单纯的情绪象征，而是一种外在的腐蚀力量。
它会逐渐扭曲少女的世界观和思维方式，使其陷入不断加深的绝望循环之中，而这种绝望又会进一步加快灵魂宝石变污浊的速度。
她们生活中遭遇的事件并不会导致这一循环，但会使这一循环进一步加剧\cite{apa2000dsm-iv-tr}。

在目前已经出现的魔法少女之中，抑郁构成了她们共同的精神底色，只是这种状态在每个人身上的表现形式并不相同。
比如杏子，初看之下似乎没有抑郁迹象，但上一集已经清楚揭示出，她实际上早已将自己从人类关系中抽离出来。
为了逃避过往经历带来的痛苦，她选择沉溺于即时的欲望满足中，以享乐来遮蔽创伤。
她的回忆也通过视觉形式强化了这种孤立感。
画面中，杏子常常独自蜷缩或瘫坐在巨大而空旷的空间里，或者作为唯一具有色彩的玩偶，被置于一片单色剪影之间。
与此同时，杏子曾经切身经历过饥饿，因此她对食物形成了强烈的情感依附，并将其视为安慰与安全感的来源。
这一点在作品中多次出现。
沙耶香即将魔女化之前，杏子向她递出零食；
后来杏子与小圆试图把沙耶香从魔女状态中救回时，她也带着食物接近小圆。
在她存放沙耶香身体的公寓中，堆积的大量食品包装也同样指向这一点。
因此，对杏子而言，食物并不仅仅是满足生理需求的事物，而是一种替代性的情感媒介。
她通过食物短暂地感受到与他人、与世界之间的联系，而这种联系正是她在意识层面不断拒绝的东西。
同时，进食也成为她自我安抚和调节痛苦情绪的方式。
由此看来，她之所以持续不断地进食，并非单纯的人物习惯或性格标签，而是暗示着她很可能患有某种进食障碍。
如果她仍然拥有人类身体，这种障碍甚至可能对她造成严重的身心伤害。

与杏子形成鲜明对比的是沙耶香，她会将自己的抑郁外化。
围绕沙耶香反复出现的水意象，也因此具有双重含义。
一方面，正如本集结尾她的魔女形态所揭示的那样，她对应的是《小美人鱼》的叙事原型。
这里所指并非是后来的迪士尼版本中那个历经艰难、最终实现身份转换的乐观童话\cite{mardoll2012little-mermaid}，而是安徒生\footnote{%
    译注：汉斯·克里斯蒂安·安徒生（1805 年 4 月 2 日—1875 年 8 月 4 日），丹麦作家、诗人。
    安徒生以其童话作品闻名于世，童话中带有人文色彩与哲学思想。
}原作中更为残酷的故事：
少女为了与所爱之人在一起而牺牲自己的人鱼之身，到头来他却全然不知她的牺牲，另择他人为伴\cite{andersen1999little-mermaid}。
另一方面，水意象也指向沙耶香自身的心理状态。
她并不是单纯地与水相连，而是仿佛正在被水吞没。
接连而来的失落、打击与幻灭，使她感到自己正在下沉。
本集开场时，她已经承受了这样的失望：
恭介从离开医院那一刻起，便基本无视了她；
她也意识到自己并没有成为理想中那个光辉而纯粹的正义化身，反而只是一名在魔法少女中并不强大的角色。
雪上加霜的是，抑郁改变了沙耶香理解世界的方式。
她不断从最负面的角度解释自身遭遇，也以最严苛的方式评价自己的行动。
因此，身体的变化在她看来意味着人性的彻底丧失，剥夺了她的全部价值；
小圆出于友情的关心，也被她误解为居高临下的怜悯；
列车上那些渣男，则被她看作是整个世界已不再值得奋力拯救的缩影。

沙耶香崩溃的关键，是她为了治愈恭介而许下愿望，却并没有因此与恭介建立起她所期待的关系。
作品没有直接呈现恭介与仁美的谈话，这一点十分重要。
尽管上一集和《叛逆》都暗示两人后来确实开始交往，但在当时，沙耶香并不可能知道这一结果。
她之所以认定事实如此，一方面是因为抑郁使她习惯于把一切都往最坏的结果想，另一方面也是因为强烈的自卑感使她不断贬低自己、抬高他人。
对沙耶香来说，仁美不可能失败。
因为在她扭曲的思维中，像沙耶香这样的人永远无法得到自己想要的东西，而其他人似乎总能顺利获得幸福\footnote{%
    译注：这里的重点并不在于事实是否已经发生，而在于沙耶香如何理解事实。
    她不是基于确凿的信息作出判断，而是基于自我否定和绝望预先得出了结论。
    列车上的爆发也应当从这一心理结构中理解。
}。
那些渣男谈论如何冷酷地抛弃真心爱自己的女性，这番话对沙耶香而言太切中自身处境，她怎么不可能不爆发？
她很可能杀死了那些男人；这一行为并不只是对这些渣男言论的惩罚，也是在间接宣泄她对恭介的愤怒。
正因如此，她随后才会说，自己已经不记得当初究竟是为了什么而战斗了。

与沙耶香的崩溃不同，本集中还呈现了另一种崩溃，那就是焰的崩溃。
她与沙耶香一样，都已在长期的精神压力下逼近极限。
事实上，在整个系列中，作品曾多次通过极短暂的闪现画面，让观众看到焰压抑在冷静外表之下的愤怒与悲伤。
有时，类似的手法也会用在其他角色身上，例如沙耶香在避雨处与小圆对话的场景。
这些画面的性质并不完全明确。
它们当然可以被理解为焰努力隐藏的真实情绪的短暂流露。
但问题在于，这些画面是否只是在叙事层外为观众提供心理线索，还是说在本集揭晓焰真实身份之后，也可以被理解为其他时间线中相同场景曾以细微差异发生过的残影。
换言之，它们既可能是叙事层外的表现性镜头，也可能暗示了时间循环中不同版本的事件。
无论采用哪一种理解，焰在射杀丘比之后的崩溃都是一个关键时刻。
她不再维持一贯的冷漠与克制，而是以近乎哀求的方式请求小圆停止自我牺牲、珍视自己的生命。
这是作品在当前这条时间线的叙事层内，第一次真正揭开焰的伪装，让我们看到她真实的痛苦、绝望和脆弱。
就连小圆也意识到了这一场面的特殊意义，因为这是她第一次清楚地意识到，自己曾经与焰相遇过。

在多数情节中，焰总是表现得冷静、理性而高效。
这种形象可以被视为“功能性抑郁”或“恶劣心境”的呈现。
与人们更熟悉的重性抑郁障碍相比，这类抑郁未必表现得同样剧烈，却往往更加持久\cite{apa2000dsm-iv-tr}。
与杏子类似，焰也在以某种方式进行自我麻痹。
不同的是，杏子借助食物和享乐来缓解痛苦，而焰所依赖的，则是对“拯救小圆”这一目标的执着。
焰并不真正怀有希望。
她不相信自己能够成功，也不相信自己将会成功，更不相信境况终会好转。
支撑她继续行动的，并不是乐观的信念，而是“还有事情非做不可”这一事实。
只要她仍在行动，仍在朝着那个无法实现的目标前进，她就可以不必正视自己的绝望，也不必回想自己在无数次循环中承受过的痛苦。
只要她不断前进、不断战斗，她就可以不去思考自己的孤立处境，也可以不去承认一个残酷的事实：
为了拯救小圆，她反而将自己推得离小圆越来越远。
因此，沙耶香对焰的判断其实相当准确：
她们是一样的，都已经在某种意义上放弃了生活。
二者的区别只在于，焰仍然拥有一个可以支撑自己的信念。
只要她能够不断说服自己，除了保护小圆之外，其他一切都无关紧要，她也就能说服自己，她本人的绝望与孤立感都无关紧要，从而使自己不至于变成魔女。
也正是在这个意义上，焰延续了她与麻美之间的对位关系。
麻美同样是在魔法少女的孤独中，依靠对某种使命的献身来维持自身。

即使是小圆本人，也表现出某种抑郁倾向，尤其体现为较低的自我价值感。
正如焰所指出的，她随时准备为他人牺牲自己，也就是说，她始终把他人置于自己之上。
这一形象与后来在其他时间线中出现的小圆相去甚远。
这似乎表明，渗入小圆意识的，不仅有其他时间线以梦境或闪回形式呈现的经验，还有成为魔法少女所内含的某种抑郁感。
因此，作品对小圆抑郁倾向的描写之所以显得不完整甚至不一致，正可以理解为她在这一时间线中尚未成为魔法少女的结果。

作品对抑郁经历的呈现既准确又多样，尤其难得的是，它并不是从外部观察抑郁，而是仿佛从抑郁者的内部视角出发，呈现其感知世界和理解自我的方式。
这种书写很容易让人联想到创作者自身可能就有相关的抑郁经历。
尤其是在结合虚渊玄于 \emph{Fate/Zero} 第一卷后记中的相关表述后——
这一点我已在第四章和第五章中讨论过——
这种联想便显得更具解释力。
不过，必须强调的是，批评者不应仅根据作品诊断艺术家的心理状态。
这样做不仅缺乏审慎，也极不尊重创作者本人。
因此，在分析中有必要严格区分“隐含作者”和“实际作者”。
前者是观众透过作品建构出来的作者形象；
后者则是真实存在的具体个人，除了真正认识他们的人之外，任何人都无法真正了解。
更何况，《魔圆》的实际作者本就是许多人，包括编剧、导演、角色设计师、动画制作者、声优等；而隐含作者则一如既往，仍是一个个体。
话虽如此，隐含的“虚渊玄”确实显得深陷抑郁与绝望之中，以至于他甚至无法想象一个积极向上、会有好事发生的世界，也无法创作描绘这样一个世界的艺术作品。

在后续几集中，丘比的真实目的逐渐显现出来：
它试图借助魔法少女与魔女产生的能量来对抗熵，从而延缓宇宙必然走向衰败的过程。
更准确地说，丘比的做法是将原本属于物理宇宙的衰败，转嫁到魔法少女的情感和精神状态之中。
于是，这种宇宙层面的衰败，在作品中就等同于抑郁和绝望。
但这个正在衰败的宇宙本身，又存在于一个深陷抑郁的隐含作者心中；
他正在奋力突破自身的局限。
换句话说，丘比所执行的机制试图以如下方式阻止并驱除抑郁：
将原本属于作者心灵的抑郁，转嫁到角色的心灵之中。

不过，本集也是其他角色开始明确把丘比视为敌人的一集，小圆除外。
尤其是在本集结尾，作品通过画面调度、音乐处理和节奏安排，明确将丘比塑造成一个彻底的恶棍。
它自鸣得意地从魔法少女的痛苦中获利，而这些少女在订立契约时，甚至并不知道自己真正付出的代价是什么。
但值得注意的是，紧接着列车上的渣男片段，便是丘比试图诱导小圆许愿拯救沙耶香。
在前者中，两位渣男谈论如何操控自己身边的女性，并将女性视为地位低于自己的对象。
沙耶香随后所说的“我已经不记得自己当初想拯救什么了”，由此把丘比的恶行嵌入这样一个问题：
人类真的比丘比更好吗？

然而，如果说丘比的职责是守护并维系隐含作者处理抑郁的方式，那么这一设定也使作者自身受到同样负面的审视。
这进一步说明，他清楚地意识到自身内在的问题，并有意将这种问题暴露出来。
换言之，隐含作者是在寻找另一种可能的出路。

然而，要找到这种出路，就必须首先瓦解他自己所建立的系统，也就是整个《魔法少女小圆》的世界结构。
要做到这一点，作品就必须先摧毁支撑它成立的框架，也即“魔法少女”类型本身。
